\documentclass[11pt,a4paper]{report}

\usepackage[margin=1in]{geometry}
\usepackage{amsmath,amssymb,amsfonts}
\usepackage{booktabs}
\usepackage{hyperref}
\usepackage{enumitem}
\usepackage{graphicx}
\usepackage{longtable}

\hypersetup{
    colorlinks=true,
    linkcolor=blue,
    citecolor=blue,
    urlcolor=blue
}

\title{OD Estimation on the Sioux Falls Network\\[0.4em]
\large Method Documentation and Machine-Learning Pipeline}
\author{Implementation: \texttt{src/}, \texttt{generate\_od\_dataset.py}, \texttt{generate\_turning\_counts.py}}
\date{\today}

% --- Notation ---
\newcommand{\Tbase}{\mathbf{T}^{\mathrm{base}}}
\newcommand{\Tsyn}{\mathbf{T}^{\mathrm{syn}}}
\newcommand{\Tsynk}[1]{\mathbf{T}^{\mathrm{syn},(#1)}}
\newcommand{\Xipf}[1]{\mathbf{X}^{(#1)}}
\newcommand{\mathbfp}{\mathbf{p}}
\newcommand{\mathbfa}{\mathbf{a}}
\newcommand{\mathbfS}{\mathbf{S}}
\newcommand{\mathbfM}{\mathbf{M}}
\newcommand{\mathbfy}{\mathbf{y}}
\newcommand{\mathbfx}{\mathbf{x}}
\newcommand{\mathbfX}{\mathbf{X}}
\newcommand{\mathbfY}{\mathbf{Y}}
\newcommand{\Aturn}{\mathbf{A}_{\mathrm{turn}}}

\begin{document}
\maketitle

\begin{abstract}
This document describes the full pipeline for origin--destination (OD) estimation on the Sioux Falls network: (i)~synthetic OD matrix generation with quality metrics, (ii)~probabilistic turning-count observation via a route-choice-aware linear assignment operator, and (iii)~machine-learning model selection for the inverse problem. The exposition is organized in three parts; a fourth part is reserved for future project stages. Notation is fixed throughout. The primary synthetic generator uses Dirichlet sampling around the base OD distribution, marginal perturbation, and iterative proportional fitting (IPF). Turning counts are generated by $K$-shortest-path enumeration, multinomial logit route choice, and a fractional assignment matrix $\Aturn$, with optional Poisson observation noise. Implementations live in \texttt{src/} and are designed for datasets of up to $10^6$ samples using batched, memory-mapped I/O.
\end{abstract}

\tableofcontents
\newpage

%=============================================================================
\part{Synthetic Data Generation and Evaluation}
\label{part:synthetic}
%=============================================================================

\chapter{Introduction and Problem Setting}
\label{ch:intro}

\section{The OD matrix}

Let $N=24$ be the number of traffic analysis zones in Sioux Falls. The \emph{base} OD matrix
\begin{equation}
    \Tbase = \bigl(T^{\mathrm{base}}_{ij}\bigr)_{i,j=1}^{N} \in \mathbb{R}_{\ge 0}^{N \times N}
\end{equation}
records daily trips from origin zone $i$ to destination zone $j$. We observe a single matrix \texttt{EstimatedODMatrix.npy}; the goal is to train models that recover OD demand from turning-count observations.

\paragraph{Marginals and grand total.}
\begin{align}
    p_i^{\mathrm{base}} &= \sum_{j=1}^{N} T^{\mathrm{base}}_{ij}
    && \text{(productions)}, \\
    a_j^{\mathrm{base}} &= \sum_{i=1}^{N} T^{\mathrm{base}}_{ij}
    && \text{(attractions)}, \\
    G &= \sum_{i,j} T^{\mathrm{base}}_{ij}
    = \sum_i p_i^{\mathrm{base}} = \sum_j a_j^{\mathrm{base}}.
\end{align}

\section{Why synthetic data?}

Machine learning requires many labeled examples. We generate $K$ synthetic matrices $\{\Tsynk{k}\}_{k=1}^{K}$ that:
\begin{itemize}[noitemsep]
    \item preserve total demand $G$ and sparsity structure;
    \item retain dominant travel patterns but explore realistic variation;
    \item pair with turning counts via a known forward model for supervised training.
\end{itemize}

\section{End-to-end pipeline overview}

\begin{figure}[ht]
\centering
\fbox{\begin{minipage}{0.94\linewidth}
\small
\textbf{Part I:} $\Tbase \to$ Dirichlet+IPF $\to$ $\{\Tsynk{k}\}_{k=1}^{K}$ (shape $K\times 24\times 24$)\\[0.3em]
\textbf{Part II:} $\Tsynk{k} \to \mathbfy^{(k)} \to \mathbfx^{(k)} = \Aturn\mathbfy^{(k)} + \boldsymbol{\epsilon}^{(k)}$ \quad ($K$-paths + logit + fractional $\Aturn$)\\[0.3em]
\textbf{Part III:} train $f:\mathbfx \mapsto \mathbfy$ on $(\mathbfX,\mathbfY)$; evaluate on held-out synthetics
\end{minipage}}
\caption{Three-part pipeline. Part~IV (future) will cover deployment, real-data validation, and equilibrium extensions.}
\label{fig:pipeline_overview}
\end{figure}

\section{Notation summary}

\begin{table}[ht]
\centering
\caption{Primary notation used throughout this document.}
\label{tab:notation}
\small
\begin{tabular}{@{}ll@{}}
\toprule
\textbf{Symbol} & \textbf{Meaning} \\
\midrule
$\Tbase$ & Observed base OD matrix \\
$\Tsyn$, $\Tsynk{k}$ & Synthetic OD (single / $k$-th replicate) \\
$\mathbfp^{\mathrm{base}}$, $\mathbfa^{\mathrm{base}}$ & Base productions and attractions \\
$\mathbfp^{\mathrm{tgt}}$, $\mathbfa^{\mathrm{tgt}}$ & IPF target marginals (sum to $G$) \\
$p_i^{\mathrm{syn}}$, $a_j^{\mathrm{syn}}$ & Achieved marginals of $\Tsyn$ \\
$\mathbfp^{\mathrm{prob}}$ & Base OD probability vector ($|\mathcal{S}|=552$ nonzero of 576) \\
$\mathbfS$, $\Xipf{t}$ & IPF seed matrix / iterate at step $t$ \\
$\mathbfM$ & Binary sparsity mask from $\Tbase$ \\
$\mathbfy$, $\mathbfY$ & Flattened OD vector / batch ($N^2=576$) \\
$\mathbfx$, $\mathbfX$ & Turning-count vector / batch ($M=178$); clean or noisy \\
$\mathbfx^{\mathrm{clean}}$, $\mathbfx^{\mathrm{obs}}$ & Clean turning counts / noisy observations \\
$\Aturn$ & Fractional turning assignment matrix $(178 \times 576)$ \\
$\Aturn^{\mathrm{bin}}$ & Binary (shortest-path) assignment matrix \\
$\mathcal{R}_{ij}$ & Set of enumerated routes for OD pair $(i,j)$ \\
$K$ & Number of shortest simple paths per OD pair \\
$\theta$ & Logit sensitivity (cost scale in route choice) \\
$c_r$, $U_r$, $P_r$ & Route cost, utility, choice probability \\
$\mathbf{b}_r$ & Binary turn-incidence vector for route $r$ \\
$\alpha$, $\delta$ & Dirichlet concentration / marginal perturbation level \\
\bottomrule
\end{tabular}
\end{table}


\chapter{Generative Methodology: From Base OD to Synthetic Matrices}
\label{ch:generation}

\section{Structural constraints}

\paragraph{Sparsity support.}
\begin{equation}
    \mathcal{S} = \{(i,j) : T^{\mathrm{base}}_{ij} > 0\}, \qquad
    M_{ij} = \mathbb{1}[(i,j)\in\mathcal{S}].
\end{equation}
All generators enforce $T^{\mathrm{syn}}_{ij}=0$ for $(i,j)\notin\mathcal{S}$.

\paragraph{Zero intrazonal trips.}
$T_{ii}^{\mathrm{syn}} = 0$ for all $i$.

\section{Two generations of synthetic generators}

The project implements two generators. Table~\ref{tab:gen_compare} compares them; the \textbf{Dirichlet+IPF} pipeline (\texttt{src/generate\_od\_dataset.py}) is the recommended default.

\begin{table}[ht]
\centering
\caption{Comparison of synthetic OD generation methods.}
\label{tab:gen_compare}
\small
\begin{tabular}{@{}p{3.2cm}p{4.8cm}p{4.8cm}@{}}
\toprule
\textbf{Criterion} & \textbf{Legacy: uniform seed noise} & \textbf{Current: Dirichlet + IPF} \\
 & (\texttt{synthetic\_od.py}) & (\texttt{generate\_od\_dataset.py}) \\
\midrule
Seed construction &
$T^{\mathrm{base}}_{ij}\cdot\eta_{ij}$, $\eta_{ij}\sim\mathcal{U}(0.8,1.2)$ &
$\tilde{\mathbfp}\sim\mathrm{Dir}(\alpha\mathbfp^{\mathrm{prob}})$, $\tilde{\mathbfT}=G\tilde{\mathbfp}$ \\
Marginal targets &
Perturb $\mathbfp^{\mathrm{base}}$, $\mathbfa^{\mathrm{base}}$; rescale to $G$ &
Same marginal perturbation (configurable $\delta$) \\
Balancing &
IPF to match targets &
IPF with Dirichlet sample as seed \\
Typical $\rho$ vs.\ base &
$\approx 0.97$ (high, limited diversity) &
$\approx 0.74$ at $\alpha=500$ ( richer ML training) \\
Pairwise diversity &
$\bar{d}\approx 139$ &
$\bar{d}\approx 545$ \\
Scalability &
Multiprocessing; 100k in $\sim$5\,s &
Multiprocessing; 50k in $\sim$2.4\,s; 1M in $\sim$1\,min \\
Hyperparameter search &
Manual &
Grid over $\alpha$, $\delta$ built-in \\
\bottomrule
\end{tabular}
\end{table}

\paragraph{Why Dirichlet+IPF was selected.}
The legacy generator produces synthetics that are almost rescaled copies of $\Tbase$ ($\rho\approx 0.97$). For ML, labels must span enough of the 576-dimensional OD space without violating transport realism. Dirichlet sampling on the base probability simplex provides a \emph{principled} way to draw alternative trip tables that remain close to $\Tbase$ when $\alpha$ is large and become more diverse when $\alpha$ is small. IPF then enforces zone-level consistency via perturbed marginals. This separates \emph{distributional} diversity (Dirichlet) from \emph{aggregate} realism (marginals + sparsity).

\section{Stage 1: Base OD analysis}

Given $\Tbase$, compute:
\begin{itemize}[noitemsep]
    \item grand total $G$;
    \item production vector $\mathbfp^{\mathrm{base}}$ and attraction vector $\mathbfa^{\mathrm{base}}$;
    \item sparsity fraction ($\approx 4.2\%$ zeros for Sioux Falls);
    \item top 10 / top 20 OD pairs by flow.
\end{itemize}
Implementation: \texttt{src/od\_analysis.py}. Outputs \texttt{base\_od\_analysis.json} and CSV summaries.

\section{Stage 2: Base probability distribution}

Flatten $\Tbase$ (row-major, C-order) to a vector of length $N^2=576$ and normalize:
\begin{equation}
    p_k = \frac{T^{\mathrm{base}}_k}{G}, \qquad \sum_{k=1}^{576} p_k = 1,
\end{equation}
with $p_k=0$ for off-support cells. Only the $|\mathcal{S}|=552$ positive entries participate in sampling. Implementation: \texttt{src/od\_probability.py}.

\section{Stage 3: Dirichlet sampling}

Draw a probability vector on the support:
\begin{equation}
    \tilde{\mathbfp} \sim \mathrm{Dirichlet}(\alpha\,\mathbfp^{\mathrm{prob}}_{\mathcal{S}}),
    \label{eq:dirichlet}
\end{equation}
where $\mathbfp^{\mathrm{prob}}_{\mathcal{S}}$ is the base probability restricted to $\mathcal{S}$, and $\alpha>0$ is a \emph{concentration} parameter.

\paragraph{Effect of $\alpha$.}
\begin{itemize}[noitemsep]
    \item \textbf{Small $\alpha$} (e.g.\ 100): higher variance, lower correlation with $\Tbase$, greater ML diversity.
    \item \textbf{Large $\alpha$} (e.g.\ 2000): samples concentrate near $\mathbfp^{\mathrm{prob}}$, similar to legacy generator.
    \item \textbf{Default $\alpha=500$}: balance between realism ($\rho\approx 0.74$) and diversity ($\bar{d}\approx 545$).
\end{itemize}

\begin{table}[ht]
\centering
\caption{Dirichlet concentration $\alpha$: trade-offs observed on Sioux Falls ($K=5000$, $\delta=0.20$).}
\label{tab:alpha}
\begin{tabular}{@{}llll@{}}
\toprule
$\alpha$ & Mean $\rho$ vs.\ base & Mean pairwise distance & Recommendation \\
\midrule
100  & lower ($\sim$0.65) & highest & exploratory / high-diversity ML \\
250  & moderate           & high    & alternative if $\alpha=500$ too conservative \\
500  & $\approx 0.74$     & $\approx 545$ & \textbf{default production setting} \\
1000 & higher             & moderate & closer to base; less label noise \\
2000 & highest ($\sim$0.90+) & lowest & ablation / sanity check only \\
\bottomrule
\end{tabular}
\end{table}

\section{Stage 4: Demand reconstruction}

Rescale sampled probabilities to trip counts and reshape to $N\times N$:
\begin{equation}
    \tilde{T}_{ij} = G \cdot \tilde{p}_{\ell(i,j)}.
\end{equation}

\section{Stage 5--6: Sparsity and diagonal enforcement}

Apply mask and zero diagonal \emph{before} IPF (seed preparation):
\begin{align}
    S_{ij} &\leftarrow \tilde{T}_{ij} \cdot M_{ij}, &
    S_{ii} &\leftarrow 0, &
    S_{ij} &\leftarrow S_{ij} + \varepsilon M_{ij}, \quad \varepsilon=10^{-6}.
\end{align}
The floor $\varepsilon$ prevents division-by-zero in IPF scaling on the support.

\section{Stage 7: Marginal perturbation}

Independent multiplicative shocks on zone totals, then rescaling to $G$:
\begin{align}
    \tilde{p}_i &= p_i^{\mathrm{base}}(1+\xi_i^p), & \xi_i^p &\sim \mathcal{U}(-\delta,\delta), \\
    \tilde{a}_j &= a_j^{\mathrm{base}}(1+\xi_j^a), & \xi_j^a &\sim \mathcal{U}(-\delta,\delta), \\
    p_i^{\mathrm{tgt}} &= G\,\tilde{p}_i \Big/ {\textstyle\sum_{i'}\tilde{p}_{i'}}, &
    a_j^{\mathrm{tgt}} &= G\,\tilde{a}_j \Big/ {\textstyle\sum_{j'}\tilde{a}_{j'}}.
\end{align}
Default $\delta=0.20$; grid search tests $\delta\in\{0.05,0.10,0.20,0.30\}$.

\paragraph{Key distinction.}
$\mathbfp^{\mathrm{tgt}}$ and $\mathbfa^{\mathrm{tgt}}$ are \emph{IPF targets}, not the final row/column sums of $\Tsyn$. IPF \emph{induces} cell-level flows from the Dirichlet seed subject to these targets.

\section{Stage 8: IPF balancing}

Given seed $\mathbfS$ and targets, IPF alternates row and column scaling (Furness/Deming--Stephan). Let $\Xipf{t}$ denote the iterate after $t$ cycles:
\begin{align}
    X^{(t+\frac{1}{2})}_{ij} &= X^{(t)}_{ij}\cdot \frac{p_i^{\mathrm{tgt}}}{\sum_{j'}X^{(t)}_{ij'}}, \\
    X^{(t+1)}_{ij} &= X^{(t+\frac{1}{2})}_{ij}\cdot \frac{a_j^{\mathrm{tgt}}}{\sum_{i'}X^{(t+\frac{1}{2})}_{i'j}}.
\end{align}
Stop when marginal error $e^{(t)}<\tau=10^{-3}$ or after $t_{\max}=1000$ iterations. Post-process: $\Tsyn=\Xipf{t}\odot\mathbfM$ with zero diagonal.

\begin{figure}[ht]
\centering
\fbox{\begin{minipage}{0.94\linewidth}
\small
\textbf{Algorithm 1:} Dirichlet+IPF synthetic OD (one replicate $k$)\\[0.4em]
\begin{enumerate}[leftmargin=*, nosep]
    \item Build $\mathbfp^{\mathrm{prob}}$ from $\Tbase$; cache mask $\mathbfM$
    \item Draw $\tilde{\mathbfp}\sim\mathrm{Dir}(\alpha\mathbfp^{\mathrm{prob}}_{\mathcal{S}})$; form seed $\mathbfS=G\,\mathrm{reshape}(\tilde{\mathbfp})$
    \item Apply sparsity, diagonal zero, floor $\varepsilon$
    \item Draw perturbed targets $(\mathbfp^{\mathrm{tgt}},\mathbfa^{\mathrm{tgt}})$ with level $\delta$; rescale to $G$
    \item $\Tsynk{k}\leftarrow\mathrm{IPF}(\mathbfS,\mathbfp^{\mathrm{tgt}},\mathbfa^{\mathrm{tgt}})$
    \item Project onto support; zero diagonal
\end{enumerate}
\end{minipage}}
\caption{Production synthetic OD generator (\texttt{src/od\_generator.py}).}
\label{alg:dirichlet_ipf}
\end{figure}


\chapter{Evaluation Metrics and Hyperparameter Selection}
\label{ch:metrics}

\section{Per-matrix metrics}

For each $\Tsyn$, with $d_{ij}=T^{\mathrm{syn}}_{ij}-T^{\mathrm{base}}_{ij}$:

\paragraph{Similarity.}
\begin{align}
    \mathrm{MAE} &= \frac{1}{N^2}\sum_{i,j}|d_{ij}|, &
    \mathrm{RMSE} &= \sqrt{\frac{1}{N^2}\sum_{i,j}d_{ij}^2}, \\
    \mathrm{RelErr} &= \frac{\sum_{i,j}|d_{ij}|}{G}, &
    \rho &= \mathrm{corr}(\mathrm{vec}(\Tbase),\mathrm{vec}(\Tsyn)).
\end{align}

\paragraph{Marginal drift from base.}
\begin{equation}
    E_p = \frac{1}{N}\sum_i\left|\frac{p_i^{\mathrm{syn}}-p_i^{\mathrm{base}}}{p_i^{\mathrm{base}}+10^{-9}}\right|,
    \quad
    E_a = \frac{1}{N}\sum_j\left|\frac{a_j^{\mathrm{syn}}-a_j^{\mathrm{base}}}{a_j^{\mathrm{base}}+10^{-9}}\right|.
\end{equation}

\paragraph{Structural validity.}
New connections $N_{\mathrm{new}}=0$, diagonal violation $D_{\mathrm{diag}}=0$, sparsity matches $\Tbase$.

\paragraph{IPF diagnostics.}
Iteration count, final error $e^*$, convergence flag.

\section{Dataset-level metrics}

For batch $\{\Tsynk{k}\}_{k=1}^{K}$:
\begin{itemize}[noitemsep]
    \item \textbf{Pairwise diversity:} sampled mean/min/max L1 distance between synthetics;
    \item \textbf{Cell variability:} per-OD mean, std, coefficient of variation;
    \item \textbf{Demand consistency:} $\mathrm{CV}_G=\mathrm{std}(\sum_{ij}T^{(k)}_{ij})/\bar{G}\approx 0$;
    \item \textbf{PCA:} explained variance; components needed for 90\%;
    \item \textbf{KMeans:} inertia for $k\in\{2,5,10\}$ to detect multimodal structure.
\end{itemize}
For $K>10^5$, metrics are computed on a random subset (default 10{,}000) for tractability.

\section{Hyperparameter grid search}

\texttt{src/hyperparameter\_search.py} evaluates all pairs $(\alpha,\delta)$ on a smaller batch (default 200 samples per config) and ranks configurations by a composite score:
\begin{equation}
    \mathrm{score} = \frac{\bar{d}}{200} + \overline{\mathrm{CV}}_{\mathrm{cell}} + \frac{n_{90\%}}{20} - 2\max(0,\bar{\rho}-0.85),
\end{equation}
rewarding diversity and PCA spread while penalizing excessive correlation with $\Tbase$. Recommended defaults from production runs: $\alpha=500$, $\delta=0.20$.

\section{Scalability optimizations}

For $K=10^6$ samples:
\begin{itemize}[noitemsep]
    \item multiprocessing (\texttt{--workers}, default all CPU cores);
    \item \texttt{float32} storage ($\sim$2.3\,GB vs.\ 4.6\,GB);
    \item per-index RNG $\mathrm{seed}+k$ for reproducibility independent of worker count;
    \item in-place IPF and masking; no IPF history stored at scale;
    \item sampled evaluation (\texttt{--eval-samples 10000}); skip plots (\texttt{--skip-plots}).
\end{itemize}
Benchmark: 50{,}000 matrices in $\approx$2.4\,s on 20 cores; 1{,}000{,}000 estimated at $\approx$45--60\,s generation.


\chapter{Implementation Reference (Part I)}
\label{ch:impl1}

\begin{table}[ht]
\centering
\caption{Part I: key modules and outputs.}
\label{tab:impl1}
\small
\begin{tabular}{@{}ll@{}}
\toprule
\textbf{Module / script} & \textbf{Role} \\
\midrule
\texttt{generate\_od\_dataset.py} & CLI orchestrator \\
\texttt{src/od\_analysis.py} & Base OD statistics \\
\texttt{src/od\_probability.py} & Probability vector on support \\
\texttt{src/dirichlet\_sampler.py} & Dirichlet draw + demand reconstruction \\
\texttt{src/marginal\_perturbation.py} & Target marginals \\
\texttt{src/ipf.py} & IPF balancing \\
\texttt{src/od\_generator.py} & Parallel batch generation \\
\texttt{src/od\_metrics.py} & Per-matrix and fast aggregate metrics \\
\texttt{src/dataset\_evaluation.py} & PCA, KMeans, diversity \\
\texttt{src/hyperparameter\_search.py} & Grid search over $(\alpha,\delta)$ \\
\midrule
\texttt{synthetic\_od.npy} & Output array shape $(K,24,24)$ \\
\texttt{dataset\_metrics.json} & Evaluation summary \\
\texttt{dataset\_summary.md} & Human-readable report \\
\bottomrule
\end{tabular}
\end{table}

\paragraph{Example command (1M samples).}
\begin{verbatim}
python generate_od_dataset.py \
  --input EstimatedODMatrix.npy \
  --samples 1000000 --alpha 500 --perturbation 0.20 \
  --skip-plots --skip-base-analysis
\end{verbatim}


%=============================================================================
\part{Turning-Count Observation and Dataset Construction}
\label{part:turning}
%=============================================================================

\chapter{Network Representation and Turning Movements}
\label{ch:network}

\section{Sioux Falls network}

The road network is parsed from \texttt{sioux-falls.net.xml} (SUMO format) into a directed graph $\mathcal{G}=(\mathcal{V},\mathcal{E})$ using NetworkX. \textbf{No traffic simulation} (SUMO, MATSim, etc.) is used at runtime---only static network topology and connection rules.

\begin{itemize}[noitemsep]
    \item $|\mathcal{V}|=24$ junctions $J_1,\ldots,J_{24}$, mapped one-to-one to zones;
    \item $|\mathcal{E}|=76$ directed regular edges with length $L_e$ and free-flow speed;
    \item 384 SUMO \texttt{<connection>} elements defining legal edge-to-edge turns.
\end{itemize}
Implementation: \texttt{src/network\_parser.py}.

\section{Turning movement definition}

A turning movement is a consecutive edge pair $(e^{\mathrm{in}},e^{\mathrm{out}})$ at junction $v$:
\begin{equation}
    \omega(e^{\mathrm{in}})=v=\alpha(e^{\mathrm{out}}), \qquad e^{\mathrm{in}}\neq e^{\mathrm{out}},
    \label{eq:turn_def}
\end{equation}
where $\omega(\cdot)$ and $\alpha(\cdot)$ denote the head and tail junction of a directed edge. Movement $m$ is \emph{admissible} if $(e^{\mathrm{in}},e^{\mathrm{out}})$ is both a valid junction turn and an allowed SUMO connection. This yields $M=178$ movements (not 384, because only physically consistent, non-U-turn connections are retained). Each movement receives a unique index $m\in\{0,\ldots,M-1\}$ stored in \texttt{turning\_id\_map}. Implementation: \texttt{src/turning\_movements.py}; catalog exported as \texttt{turning\_movements.csv}.

\section{Two generations of route assignment}

The project implements two forward models from OD to turning counts. Table~\ref{tab:route_gen_compare} compares them; the \textbf{$K$-shortest-path + logit} pipeline (\texttt{generate\_turning\_counts.py}) is the recommended default for machine-learning dataset generation.

\begin{table}[ht]
\centering
\caption{Comparison of route assignment and turning-count generation methods.}
\label{tab:route_gen_compare}
\small
\begin{tabular}{@{}p{3.2cm}p{4.8cm}p{4.8cm}@{}}
\toprule
\textbf{Criterion} & \textbf{Phase~1: binary shortest path} & \textbf{Phase~2: probabilistic (current)} \\
 & (\texttt{generate\_dataset.py}) & (\texttt{generate\_turning\_counts.py}) \\
\midrule
Route set &
Single shortest path $\pi_{ij}$ per OD &
Up to $K$ simple shortest paths $\mathcal{R}_{ij}$ \\
Route choice &
Deterministic (all-or-nothing) &
Multinomial logit on path costs \\
$\Aturn$ entries &
Binary: $A_{m,\ell}\in\{0,1\}$ &
Fractional: $A_{m,\ell}\in[0,1]$ \\
$\mathrm{rank}(\Aturn)$ &
122 (nullity 454) &
174 (nullity 402) \\
Route diversity &
None (same path always) &
Controlled by $(K,\theta)$ \\
Observation noise &
None &
None / Gaussian / Poisson \\
Scalability at $K_{\mathrm{samples}}=10^6$ &
Full in-memory load &
Batched memory-mapped I/O \\
Typical use &
Baseline / ablation &
\textbf{Production ML labels} \\
\bottomrule
\end{tabular}
\end{table}

\paragraph{Why Phase~2 was selected.}
Phase~1 assigns every OD pair to a single shortest path. Real drivers split across feasible alternatives; a binary $\Aturn$ therefore (i)~understates route diversity, (ii)~lowers the rank and information content of the forward operator, and (iii)~produces turning-count batches with artificially high pairwise correlation. Phase~2 retains the \emph{linear} forward model $\mathbfx=\Aturn\mathbfy$---essential for scalable supervised learning---while replacing deterministic routing with a classical random utility model (multinomial logit) over a finite route set. The result is a fractional $\Aturn$ that is still matrix--vector multiplyable at $10^6$ samples but encodes richer route-choice structure.


\chapter{Route Enumeration and Probabilistic Assignment}
\label{ch:routes}

\section{OD pair indexing}

All $N^2=576$ zone pairs are indexed in row-major (C) order:
\begin{equation}
    \ell(i,j)=(i-1)N+(j-1), \qquad y_{\ell(i,j)}=T_{ij}, \qquad i,j\in\{1,\ldots,N\}.
    \label{eq:od_index}
\end{equation}
The flattened OD vector is $\mathbfy=(y_0,\ldots,y_{575})^\top\in\mathbb{R}_{\ge 0}^{576}$. Implementation: \texttt{src/od\_pairs.py}.

\section{Route catalog: $K$ shortest simple paths}

For each ordered zone pair $(i,j)$ with $i\neq j$, let $s_i$ and $t_j$ denote the corresponding junction nodes in $\mathcal{G}$. Define the \emph{route catalog}
\begin{equation}
    \mathcal{R}_{ij} = \bigl\{ r_1,\ldots,r_{K_{ij}} \bigr\}, \qquad K_{ij}\le K,
\end{equation}
where each route $r$ is an edge sequence $\pi=(e_1,\ldots,e_L)$ obtained from \texttt{networkx.shortest\_simple\_paths}, which yields simple paths in non-decreasing cost order. Edge cost is either length $L_e$ (default) or free-flow travel time $T_e=L_e/v_e$.

The cost of route $r$ is
\begin{equation}
    c_r = \sum_{e\in r} w_e, \qquad w_e \in \{L_e, T_e\}.
    \label{eq:route_cost}
\end{equation}
Default $K=5$; sensitivity analysis tests $K\in\{1,3,5,10\}$. Implementation: \texttt{src/k\_shortest\_paths.py}; full catalog saved as \texttt{route\_catalog.pkl}.

\begin{table}[ht]
\centering
\caption{Route enumeration methods: pros and cons.}
\label{tab:route_enum}
\small
\begin{tabular}{@{}p{3.5cm}p{4.5cm}p{4.5cm}@{}}
\toprule
\textbf{Method} & \textbf{Pros} & \textbf{Cons} \\
\midrule
Single shortest path &
Fast; deterministic baseline &
Zero route diversity \\
$K$ shortest \emph{simple} paths (chosen) &
Captures discrete alternatives; no loops; cost-ordered &
$K$ hyperparameter; enumeration cost grows with network \\
$K$ shortest \emph{with} loops &
More paths &
Unrealistic cycling; excluded \\
Random walk sampling &
Cheap &
No cost ordering; hard to interpret \\
Simulation-based routing &
Realistic congestion &
Breaks linear model; slow at $10^6$ labels \\
\bottomrule
\end{tabular}
\end{table}

\paragraph{Why $K$ shortest simple paths.}
Yen's / NetworkX simple-path enumeration provides a \emph{simulation-free}, reproducible set of cost-ranked alternatives. Unlike user equilibrium, it preserves a fixed route catalog per OD pair, so $\Aturn$ remains constant across the dataset and the forward map stays linear in $\mathbfy$.

\section{Multinomial logit route choice}

Given $\mathcal{R}_{ij}$ with costs $\{c_r\}$, assign utility and probability
\begin{align}
    U_r &= -\theta\, c_r, \label{eq:utility}\\
    P_r &= \frac{\exp(U_r)}{\sum_{r'\in\mathcal{R}_{ij}} \exp(U_{r'})}
    = \frac{\exp(-\theta c_r)}{\sum_{r'\in\mathcal{R}_{ij}} \exp(-\theta c_{r'})},
    \label{eq:logit}
\end{align}
where $\theta>0$ is a \emph{sensitivity} parameter. Probabilities are computed via a numerically stable softmax (subtract $\max_r U_r$ before exponentiation). Default $\theta=0.1$; sensitivity grid $\theta\in\{0.05,0.1,0.2,0.5\}$.

\paragraph{Interpretation of $\theta$.}
\begin{itemize}[noitemsep]
    \item \textbf{Small $\theta$} (e.g.\ 0.05): utilities are flat; routes split more evenly $\Rightarrow$ higher route diversity and higher $\mathrm{rank}(\Aturn)$.
    \item \textbf{Large $\theta$} (e.g.\ 0.5): probability mass concentrates on minimum-cost routes $\Rightarrow$ Phase~2 collapses toward Phase~1.
\end{itemize}
Implementation: \texttt{src/route\_choice.py}.

\begin{table}[ht]
\centering
\caption{Route choice models: pros and cons.}
\label{tab:route_choice}
\small
\begin{tabular}{@{}p{3.5cm}p{4.5cm}p{4.5cm}@{}}
\toprule
\textbf{Model} & \textbf{Pros} & \textbf{Cons} \\
\midrule
All-or-nothing (Phase~1) &
Trivial; rank-122 baseline &
Unrealistic; low ML signal \\
Multinomial logit (chosen) &
Standard random utility; closed form; one parameter &
IIA assumption; no congestion feedback \\
Nested logit &
Captures route correlation &
More parameters; heavier calibration \\
Mixed logit &
Random taste heterogeneity &
Simulation required for choice probs \\
User equilibrium &
Network-consistent flows &
Nonlinear; not fixed $\Aturn$ \\
Proportional to $1/c_r$ &
Simple &
No random utility foundation \\
\bottomrule
\end{tabular}
\end{table}

\paragraph{Why multinomial logit.}
Logit is the canonical closed-form discrete choice model (Ort\'{u}zar and Willumsen). With a single scale parameter $\theta$ it interpolates between uniform route splitting ($\theta\to 0$) and deterministic shortest-path assignment ($\theta\to\infty$), enabling controlled sensitivity analysis without re-running a traffic simulator.


\chapter{Fractional Assignment Matrix and Forward Model}
\label{ch:assignment}

\section{Route-level turn incidence}

For route $r=(e_1,\ldots,e_L)$, define consecutive turns $\Theta_r=\{(e_k,e_{k+1})\}_{k=1}^{L-1}$. The binary turn-incidence vector $\mathbf{b}_r\in\{0,1\}^M$ has
\begin{equation}
    b_{r,m} = \mathbb{1}\bigl[m \text{ corresponds to some turn in } \Theta_r\bigr].
    \label{eq:route_turn_vector}
\end{equation}
Implementation: \texttt{src/fractional\_assignment.py}.

\section{Fractional turning assignment matrix}

Let $\Aturn\in[0,1]^{M\times N^2}$ denote the \emph{probabilistic} assignment matrix. For OD index $\ell=\ell(i,j)$:
\begin{equation}
    A_{\mathrm{turn},m,\ell}
    = \sum_{r\in\mathcal{R}_{ij}} P_r\, b_{r,m}.
    \label{eq:fractional_A}
\end{equation}
Column $\ell$ sums to the total probability mass on routes for pair $(i,j)$ (unity when routes exist). Entries lie in $[0,1]$; shape $(178,576)$ with 7{,}326 nonzero entries (vs.\ 1{,}324 in Phase~1). Implementation: \texttt{src/fractional\_assignment.py}; saved as \texttt{A\_turn.npy}.

\paragraph{Relation to Phase~1.}
The binary matrix $\Aturn^{\mathrm{bin}}$ satisfies
\begin{equation}
    A^{\mathrm{bin}}_{m,\ell} = \mathbb{1}\bigl[m \in \Theta_{\pi_{ij}}\bigr],
\end{equation}
where $\pi_{ij}$ is the unique shortest path. Phase~2 generalizes this by weighting each route by $P_r$.

\section{Forward model (clean turning counts)}

For one synthetic OD matrix $\Tsyn$ with flattened demand $\mathbfy$:
\begin{equation}
    \mathbfx^{\mathrm{clean}} = \Aturn\,\mathbfy,
    \qquad
    x^{\mathrm{clean}}_m = \sum_{\ell=0}^{575} A_{\mathrm{turn},m,\ell}\, y_\ell.
    \label{eq:forward_single}
\end{equation}
For a batch of $K_{\mathrm{samples}}$ replicates:
\begin{equation}
    \mathbfX^{\mathrm{clean}} = \mathbfY\,\Aturn^\top,
    \qquad
    \mathbfY\in\mathbb{R}^{K_{\mathrm{samples}}\times 576},\;
    \mathbfX^{\mathrm{clean}}\in\mathbb{R}^{K_{\mathrm{samples}}\times M}.
    \label{eq:forward_batch}
\end{equation}
Equation~\eqref{eq:forward_batch} is identical in form to Phase~1; only $\Aturn$ changes from binary to fractional. Turning count $x_m$ aggregates OD flows weighted by the probability that each route uses movement~$m$.

\section{Identifiability and rank analysis}

Linear-algebra diagnostics compare Phase~1 and Phase~2 operators:

\begin{table}[ht]
\centering
\caption{Assignment matrix rank comparison (Sioux Falls, $K=5$, $\theta=0.1$).}
\label{tab:rank_compare}
\begin{tabular}{@{}lcc@{}}
\toprule
\textbf{Quantity} & $\Aturn^{\mathrm{bin}}$ (Phase~1) & $\Aturn$ (Phase~2) \\
\midrule
Shape & $(178,576)$ & $(178,576)$ \\
Nonzeros & 1{,}324 & 7{,}326 \\
$\mathrm{rank}(\Aturn)$ & 122 & \textbf{174} \\
Nullity ($576-\mathrm{rank}$) & 454 & 402 \\
$\kappa(\Aturn)$ & $\approx 14.8$ & $\approx 1.2\times 10^{12}$ \\
\bottomrule
\end{tabular}
\end{table}

\begin{align}
    \mathrm{rank}(\Aturn^{\mathrm{bin}}) &= 122, &
    \mathcal{N}_{\mathrm{bin}} &= 576-122 = 454, \label{eq:rank_bin}\\
    \mathrm{rank}(\Aturn) &= 174, &
    \mathcal{N} &= 576-174 = 402. \label{eq:rank_prob}
\end{align}

Rank increases by 52: probabilistic splitting adds 52 independent column directions in $\Aturn$, improving the information content of turning-count observations for OD recovery. The condition number rises because fractional columns are nearly collinear for some OD pairs---regularization in Part~III becomes even more important. Implementation: \texttt{src/assignment\_rank.py}, \texttt{src/rank\_analysis.py}; singular-value plots in \texttt{figures/singular\_values.png}.

\paragraph{Implication for Part III.}
The inverse problem remains underdetermined ($\mathcal{N}=402>0$), but the identifiable subspace is larger. Ridge, Elastic Net, and forward-consistency losses should be evaluated on Phase~2 labels to measure whether the rank gain translates to lower OD recovery error.


\chapter{Observation Noise and Dataset Assembly}
\label{ch:dataset}

\section{Observation noise models}

Real loop detectors and camera counts are noisy. Phase~2 supports three observation models applied \emph{after} clean synthesis:

\paragraph{No noise.}
$\mathbfx^{\mathrm{obs}} = \mathbfx^{\mathrm{clean}}$.

\paragraph{Gaussian noise.}
\begin{equation}
    x^{\mathrm{obs}}_m = x^{\mathrm{clean}}_m + \epsilon_m,
    \qquad \epsilon_m \sim \mathcal{N}\bigl(0,\sigma^2\,(x^{\mathrm{clean}}_m)_+\bigr),
    \label{eq:gauss_noise}
\end{equation}
with clipping at zero. Parameter $\sigma$ controls relative noise level.

\paragraph{Poisson noise (default).}
\begin{equation}
    x^{\mathrm{obs}}_m \sim \mathrm{Poisson}\bigl(x^{\mathrm{clean}}_m\bigr),
    \label{eq:poisson_noise}
\end{equation}
which respects nonnegativity and heteroscedastic variance $\mathrm{Var}(x^{\mathrm{obs}}_m\mid\mathbfx^{\mathrm{clean}})=x^{\mathrm{clean}}_m$. Nominal noise levels 0\%, 2\%, 5\%, 10\% map to a Poisson scaling parameter when partial noise is desired. Implementation: \texttt{src/observation\_noise.py}.

\begin{table}[ht]
\centering
\caption{Observation noise models: pros and cons.}
\label{tab:noise}
\small
\begin{tabular}{@{}p{3.5cm}p{4.5cm}p{4.5cm}@{}}
\toprule
\textbf{Model} & \textbf{Pros} & \textbf{Cons} \\
\midrule
None &
Exact labels for sanity checks &
Optimistic ML performance \\
Gaussian &
Simple; additive &
Can produce negative counts without clipping \\
Poisson (chosen) &
Standard for count data; mean=variance &
Not suitable for overdispersed sensors without extension \\
Negative binomial &
Handles overdispersion &
Extra dispersion parameter \\
\bottomrule
\end{tabular}
\end{table}

\paragraph{Why Poisson.}
Turning counts are non-negative integer aggregates of trips. Poisson noise is the canonical count-data model and matches the signal-to-noise structure of traffic sensors at moderate volumes (Cascetta). Gaussian noise is retained for ablation.

\section{End-to-end Phase~2 pipeline}

\begin{figure}[ht]
\centering
\fbox{\begin{minipage}{0.94\linewidth}
\small
\textbf{Network:} \texttt{sioux-falls.net.xml} $\to$ $\mathcal{G}$, 178 turns\\[0.3em]
\textbf{Route catalog:} for each $(i,j)$: $\mathcal{R}_{ij}\leftarrow$ $K$ shortest simple paths\\[0.3em]
\textbf{Logit:} $P_r \propto \exp(-\theta c_r)$\\[0.3em]
\textbf{Fractional $\Aturn$:} $A_{m,\ell}=\sum_r P_r b_{r,m}$\\[0.3em]
\textbf{Labels:} $\{\Tsynk{k}\}_{k=1}^{K_{\mathrm{samples}}}
\xrightarrow{\;\mathrm{vec}\;}\mathbfY
\xrightarrow{\;\mathbfX=\mathbfY\Aturn^\top\;}\mathbfX^{\mathrm{clean}}
\xrightarrow{\;\text{Poisson}\;}\mathbfX^{\mathrm{obs}}$
\end{minipage}}
\caption{Phase~2 turning-count dataset pipeline (\texttt{generate\_turning\_counts.py}).}
\label{fig:phase2_pipeline}
\end{figure}

\begin{figure}[ht]
\centering
\fbox{\begin{minipage}{0.94\linewidth}
\small
\textbf{Algorithm 2:} Batched turning-count generation (scale $K_{\mathrm{samples}}=10^6$)\\[0.4em]
\begin{enumerate}[leftmargin=*, nosep]
    \item Parse network; enumerate 178 turns; build $\Aturn$ (once)
    \item Memory-map \texttt{synthetic\_od.npy} (read-only); never load full $\mathbfY$
    \item For each batch $b$ of size $B$ (default $B=10{,}000$--$20{,}000$):
    \begin{enumerate}[nosep]
        \item $\mathbfY^{(b)} \leftarrow \mathrm{reshape}(\Tsyn[b:b+B])$ as \texttt{float32}
        \item $\mathbfX^{\mathrm{clean},(b)} \leftarrow \mathbfY^{(b)}\Aturn^\top$
        \item $\mathbfX^{\mathrm{obs},(b)} \leftarrow \mathrm{Poisson}(\mathbfX^{\mathrm{clean},(b)})$
        \item Write rows to memory-mapped \texttt{turning\_counts.npy}
    \end{enumerate}
    \item Streaming mean/std; sampled pairwise diversity; optional figures from aggregates
\end{enumerate}
\end{minipage}}
\caption{Memory-safe batching (\texttt{src/batch\_turning.py}). Peak RAM $\approx B(576+M)\times 4$ bytes $\approx 60$\,MB at $B=20{,}000$.}
\label{alg:batch_turning}
\end{figure}

\section{Scalability: why the system restarted and how it was fixed}

The initial Phase~2 implementation loaded all $K_{\mathrm{samples}}=10^6$ OD matrices into RAM and materialized multiple full copies ($\mathbfY$, $\mathbfX^{\mathrm{clean}}$, $\mathbfX^{\mathrm{obs}}$), with peak usage exceeding 12--15\,GB on a 15\,GB machine---triggering the Linux OOM killer.

The production fix (\texttt{src/batch\_turning.py}) uses:
\begin{itemize}[noitemsep]
    \item \textbf{Memory-mapped input:} \texttt{np.load(..., mmap\_mode='r')} for \texttt{synthetic\_od\_1m.npy} ($\sim$2.3\,GB on disk, not in RAM);
    \item \textbf{Batched matrix multiply:} $\mathbfX=\mathbfY\Aturn^\top$ in chunks of size $B$;
    \item \textbf{Memory-mapped output:} incremental write to \texttt{turning\_counts.npy} ($\sim$680\,MB);
    \item \textbf{Streaming statistics:} online mean/variance; sampled pairwise distances;
    \item \textbf{Auto-batching:} enabled for $K_{\mathrm{samples}}\ge 50{,}000$; \texttt{--batch-size} configurable.
\end{itemize}
Benchmark: $K_{\mathrm{samples}}=10^6$, $B=20{,}000$, $K=5$, $\theta=0.1$ completes in $\approx$16\,s with peak batch RAM $\approx 60$\,MB.

\section{Outputs and deliverables}

\begin{table}[ht]
\centering
\caption{Part II deliverables (\texttt{outputs/turning\_counts\_1m/}).}
\label{tab:outputs2}
\begin{tabular}{@{}ll@{}}
\toprule
\textbf{File} & \textbf{Content} \\
\midrule
\texttt{A\_turn.npy} & Fractional assignment matrix $(178,576)$, \texttt{float32} \\
\texttt{turning\_counts.npy} & $\mathbfX^{\mathrm{clean}}$, shape $(K_{\mathrm{samples}},178)$ \\
\texttt{turning\_counts\_noisy.npy} & $\mathbfX^{\mathrm{obs}}$ (Poisson default) \\
\texttt{route\_catalog.pkl} & $K$-path catalog + logit probabilities \\
\texttt{turning\_movements.csv} & Movement index $(e^{\mathrm{in}},e^{\mathrm{out}},m)$ \\
\texttt{assignment\_analysis.json} & Rank comparison, validation, timings \\
\texttt{figures/} & Histograms, heatmap, singular values, rank bar chart \\
\texttt{turning\_dataset\_summary.md} & Human-readable report \\
\bottomrule
\end{tabular}
\end{table}

\section{Validation}

\paragraph{Forward consistency.}
For each batch, verify $\|\mathbfx^{\mathrm{clean}}-\Aturn\mathbfy\|_\infty<10^{-3}$ on a probe sample.

\paragraph{Structural checks.}
178 turns; 576 OD columns; fractional entries in $[0,1]$; shapes $(K_{\mathrm{samples}},178)$ and $(K_{\mathrm{samples}},576)$ for targets loaded separately from Part~I.

\section{Sensitivity analysis}

Grid over $(K,\theta)$ measures $\mathrm{rank}(\Aturn)$, nullity, $\kappa(\Aturn)$, and maximum column entry. Recommended production defaults from Sioux Falls runs: $K=5$, $\theta=0.1$. Increasing $K$ at fixed $\theta$ adds routes but with diminishing rank gains beyond $K=5$; increasing $\theta$ moves Phase~2 toward Phase~1. Enable via \texttt{--run-sensitivity} (optional; re-enumerates paths per $K$).

\section{Part II evaluation highlights}

On $K_{\mathrm{samples}}=10^6$ synthetics (Dirichlet+IPF Part~I, Phase~2 Part~II):
\begin{itemize}[noitemsep]
    \item $\mathrm{rank}(\Aturn)=174$ vs.\ Phase~1 rank 122 ($+52$);
    \item mean turning count per movement $\approx 3{,}051$ (batch mean over first 1000 samples);
    \item OD pairwise diversity (sampled) $\bar{d}_{\mathrm{OD}}\approx 545$;
    \item turning pairwise diversity $\bar{d}_{\mathrm{turn}}\approx 894$;
    \item mean inter-turn correlation (50-feature sample) $\approx 0.01$.
\end{itemize}
Turning counts remain a linear compression of OD demand, but probabilistic assignment \emph{increases} turning-space separation relative to Phase~1 while preserving the exact forward map needed for supervised training.


\chapter{Implementation Reference (Part II)}
\label{ch:impl2}

\begin{table}[ht]
\centering
\caption{Part II: key modules and scripts.}
\label{tab:impl2}
\small
\begin{tabular}{@{}ll@{}}
\toprule
\textbf{Module / script} & \textbf{Role} \\
\midrule
\texttt{generate\_turning\_counts.py} & CLI orchestrator (Phase~2) \\
\texttt{src/network\_parser.py} & SUMO $\to$ NetworkX graph \\
\texttt{src/turning\_movements.py} & 178 movements + CSV export \\
\texttt{src/k\_shortest\_paths.py} & $K$ shortest simple paths per OD \\
\texttt{src/route\_choice.py} & Multinomial logit probabilities \\
\texttt{src/fractional\_assignment.py} & Route turn vectors + fractional $\Aturn$ \\
\texttt{src/assignment\_rank.py} & Binary vs.\ probabilistic rank analysis \\
\texttt{src/observation\_noise.py} & None / Gaussian / Poisson \\
\texttt{src/batch\_turning.py} & Memory-mapped batched generation \\
\texttt{src/turning\_statistics.py} & Stats, figures, sensitivity grid \\
\midrule
\texttt{generate\_dataset.py} & Phase~1 baseline (binary $\Aturn$) \\
\texttt{src/route\_assignment.py} & Shortest-path routes \\
\texttt{src/build\_assignment\_matrix.py} & Binary $\Aturn^{\mathrm{bin}}$ \\
\bottomrule
\end{tabular}
\end{table}

\paragraph{Example command (1M samples).}
\begin{verbatim}
python generate_turning_counts.py \
  --network sioux-falls.net.xml \
  --od outputs/od_generator/synthetic_od_1m.npy \
  --output-dir outputs/turning_counts_1m \
  --k_paths 5 --theta 0.1 --noise poisson \
  --batch-size 20000 --skip-plots
\end{verbatim}

\paragraph{Legacy Phase~1 command (binary baseline).}
\begin{verbatim}
python run_dataset_pipeline.py \
  --od synthetic_od_100000.npy --output-dir outputs/
\end{verbatim}


%=============================================================================
\part{Machine Learning Model Selection and Training}
\label{part:ml}
%=============================================================================

\chapter{Inverse Problem Formulation}
\label{ch:inverse}

\section{Problem statement}

Given turning-count observation $\mathbfx^{\mathrm{obs}}\in\mathbb{R}^{178}$, estimate OD vector $\mathbfy\in\mathbb{R}^{576}_{\ge 0}$ such that
\begin{equation}
    \mathbfx^{\mathrm{obs}} \approx \Aturn\mathbfy,
\end{equation}
where $\Aturn$ is the \emph{fractional} Phase~2 assignment matrix ($\mathrm{rank}=174$). Because $\Aturn$ is rank-deficient ($\mathcal{N}=402$), the problem is \emph{underdetermined}. Training uses synthetic pairs $(\mathbfx^{(k)},\mathbfy^{(k)})$ where $\mathbfx^{\mathrm{clean},(k)}=\Aturn\mathbfy^{(k)}$ exactly and $\mathbfx^{\mathrm{obs},(k)}$ may include Poisson noise.

\section{Learning paradigms considered}

\begin{table}[ht]
\centering
\caption{ML paradigms for OD estimation from turning counts.}
\label{tab:paradigms}
\small
\begin{tabular}{@{}p{3.5cm}p{4.5cm}p{4.5cm}@{}}
\toprule
\textbf{Paradigm} & \textbf{Pros} & \textbf{Cons} \\
\midrule
Linear regression + regularization &
Fast; interpretable; encodes $\Aturn$ structure via priors &
Limited nonlinear capacity \\
Physics-informed (constrained LS) &
Enforces $\mathbfx=\Aturn\mathbfy$; minimal black-box &
Needs strong priors; null-space ambiguity \\
Deep neural network &
Flexible; scales to large $K$ &
Risk of overfitting; may violate nonnegativity/sparsity \\
Gradient boosting &
Strong tabular performance; handles interactions &
Less interpretable; slower at inference \\
Bayesian inversion &
Uncertainty quantification &
Computationally heavier \\
\bottomrule
\end{tabular}
\end{table}


\chapter{Candidate Models and Selection Rationale}
\label{ch:models}

\section{Model zoo}

Table~\ref{tab:models} lists candidate models evaluated for this project. The \textbf{recommended starting stack} is Ridge regression (baseline) followed by a regularized MLP if nonlinear effects are needed.

\begin{table}[ht]
\centering
\caption{Candidate ML models: pros, cons, and suitability for turning-count $\to$ OD inversion.}
\label{tab:models}
\footnotesize
\begin{tabular}{@{}p{2.8cm}p{3.8cm}p{3.8cm}p{2.8cm}@{}}
\toprule
\textbf{Model} & \textbf{Pros} & \textbf{Cons} & \textbf{Selected?} \\
\midrule
Ordinary least squares &
Closed form; baseline &
Ill-posed ($M<N^2$); no priors &
No \\
\midrule
Ridge regression &
Stable under collinearity; analytic solution; fast &
Linear only &
\textbf{Yes (baseline)} \\
\midrule
Lasso / Elastic Net &
Promotes sparsity in $\mathbfy$; matches OD structure &
Linear; tuning $\lambda$ &
\textbf{Yes (sparse OD)} \\
\midrule
Non-negative LS with $\|\mathbfy\|_2,\|\mathbfy\|_1$ &
Enforces $y_\ell\ge 0$ &
Still linear; null-space &
Yes (extension) \\
\midrule
Tikhonov on $\Aturn$ &
Explicit use of forward operator &
Requires choosing penalty on null-space &
Yes (physics baseline) \\
\midrule
MLP (feedforward NN) &
Nonlinear; scales to large $K$ &
Needs NN constraints for $\ge 0$; tuning &
\textbf{Yes (nonlinear)} \\
\midrule
Random Forest / XGBoost &
Strong on tabular data; robust &
178-dim input ok but 576-dim output needs multi-output; slow &
Secondary \\
\midrule
CNN on $24\times 24$ &
Exploits grid structure if reshaped &
Turning counts lose spatial grid; awkward &
No \\
\midrule
GNN on network &
Natural for transport topology &
Heavy implementation; Part~IV &
Future \\
\bottomrule
\end{tabular}
\end{table}

\section{Why Ridge + Elastic Net as primary baselines}

\begin{enumerate}[leftmargin=*]
    \item \textbf{Ill-posedness:} with nullity $\mathcal{N}=402$ under Phase~2 (454 under Phase~1), unregularized inversion is unstable. Ridge ($L_2$) shrinks coefficients in the identifiable subspace.
    \item \textbf{Sparsity:} Elastic Net adds $L_1$ penalty, encouraging zero flows on off-support cells---consistent with $\Tbase$ structure.
    \item \textbf{Speed:} analytic / convex optimization trains in seconds even for $K=10^6$.
    \item \textbf{Interpretability:} coefficients relate directly to turning-movement sensitivities.
\end{enumerate}

\section{Why MLP as nonlinear extension}

Turning counts are \emph{linear} in OD under both Phase~1 and Phase~2 forward models, so a linear model should perform well on synthetics generated with the same $\Aturn$. An MLP is retained because:
\begin{itemize}[noitemsep]
    \item real data may include measurement noise and model mismatch;
    \item nonlinear calibration (e.g.\ log-transform of counts) may help;
    \item it stress-tests whether the synthetic dataset spans enough variation for nonlinear methods to add value over Ridge.
\end{itemize}
Constraints: ReLU + nonnegative output layer (or softplus); optional sparsity penalty on outputs.

\section{Why not simulation-based or equilibrium models for training labels?}

Traffic simulators produce realistic flows but are slow, stochastic, and couple OD with route choice nonlinearly. For generating $10^6$ labeled pairs, the linear $\Aturn$ model is \emph{sufficient and necessary} for throughput and for maintaining a known ground-truth Jacobian structure.


\chapter{Training Protocol and Evaluation}
\label{ch:training}

\section{Recommended data split}

\begin{itemize}[noitemsep]
    \item Generate $K_{\mathrm{samples}}=10^6$ synthetics (Part~I) and corresponding $(\mathbfX^{\mathrm{obs}},\mathbfY)$ (Part~II, Phase~2);
    \item Split 80/10/10 train/validation/test by sample index (seed-controlled);
    \item \textbf{Never} mix base OD as a test sample---it is not in the synthetic distribution.
\end{itemize}

\section{Feature and target preprocessing}

\begin{itemize}[noitemsep]
    \item Inputs $\mathbfX$: standardize turning counts (178 features) using training-set mean/std;
    \item Targets $\mathbfY$: optional $\log(1+y)$ transform; clip negatives to zero after prediction;
    \item Post-process: project predictions onto sparsity mask $\mathbfM$ and zero diagonal.
\end{itemize}

\section{Loss functions}

\begin{table}[ht]
\centering
\caption{Loss functions for OD recovery.}
\label{tab:loss}
\begin{tabular}{@{}lll@{}}
\toprule
\textbf{Loss} & \textbf{Formula (single sample)} & \textbf{Use case} \\
\midrule
MSE & $\|\hat{\mathbfy}-\mathbfy\|_2^2$ & Default \\
MAE & $\|\hat{\mathbfy}-\mathbfy\|_1$ & Robust to outliers \\
Forward-consistency & $\|\mathbfx-\Aturn\hat{\mathbfy}\|_2^2$ & Physics-informed fine-tuning \\
Combined & MSE $+\lambda\|\mathbfx-\Aturn\hat{\mathbfy}\|_2^2$ & Enforce observations \\
\bottomrule
\end{tabular}
\end{table}

For ill-posed inversion, \textbf{combined loss} (supervised MSE + forward-consistency) is recommended when deploying beyond pure synthetic evaluation.

\section{ML evaluation metrics}

On held-out test set:
\begin{align}
    \mathrm{RMSE}_{\mathrm{OD}} &= \sqrt{\frac{1}{N^2}\|\hat{\mathbfy}-\mathbfy\|_2^2}, \\
    \rho_{\mathrm{OD}} &= \mathrm{corr}(\mathbfy,\hat{\mathbfy}), \\
    \mathrm{RMSE}_{\mathrm{turn}} &= \|\mathbfx-\Aturn\hat{\mathbfy}\|_2, \\
    \mathrm{MAE}_{\mathrm{marginal}} &= \text{production/attraction error vs.\ ground truth}.
\end{align}

\section{Hyperparameter tuning (Part III)}

\begin{itemize}[noitemsep]
    \item Ridge / Elastic Net: cross-validate $\lambda$ on log-spaced grid;
    \item MLP: hidden layers $\{128,256\}$, dropout 0.1--0.3, Adam, early stopping on validation MSE;
    \item Compare against \textbf{Tikhonov baseline:} $\hat{\mathbfy}=(\Aturn^\top\Aturn+\lambda I)^{-1}\Aturn^\top\mathbfx$.
\end{itemize}

\section{Planned implementation (Part III)}

The following modules are reserved for the next development stage:
\begin{verbatim}
src/ml/
  datasets.py      # load X, Y; train/val/test split
  baselines.py     # Ridge, Elastic Net, Tikhonov
  mlp_model.py     # PyTorch or sklearn MLP
  train.py         # training loop + logging
  evaluate.py      # RMSE, correlation, forward error
\end{verbatim}


%=============================================================================
\part{Future Work and Extensions}
\label{part:future}
%=============================================================================

\chapter{Roadmap for Subsequent Chapters}
\label{ch:future}

This part is intentionally brief. As the project progresses, the following chapters will be added:

\begin{table}[ht]
\centering
\caption{Planned future chapters (Part IV+).}
\label{tab:future}
\begin{tabular}{@{}ll@{}}
\toprule
\textbf{Chapter (planned)} & \textbf{Topic} \\
\midrule
Real-data validation & Compare predictions against held-out counts / partial OD \\
Overdispersed noise & Negative binomial / compound Poisson on $\mathbfx^{\mathrm{obs}}$ \\
User equilibrium extension & Nonlinear forward model; simulation labels \\
Uncertainty quantification & Bayesian or ensemble CIs on $\mathbfy$ \\
Deployment & Inference pipeline from field turning counts \\
\bottomrule
\end{tabular}
\end{table}

\paragraph{Completed since initial roadmap.}
Multi-path fractional assignment ($K$-shortest + logit) and Poisson observation noise are now implemented in Part~II (\texttt{generate\_turning\_counts.py}).

\paragraph{Remaining limitations.}
\begin{itemize}[leftmargin=*]
    \item Single observed $\Tbase$; no temporal dynamics;
    \item No spatial correlation in Dirichlet shocks;
    \item Logit route choice without congestion feedback (static costs);
    \item Nullity $\mathcal{N}=402$ still limits unique OD recovery without priors;
    \item Part III training code not yet implemented (documented design only).
\end{itemize}


%=============================================================================
\chapter*{Conclusion}
\addcontentsline{toc}{chapter}{Conclusion}

Part~I produces diverse, structurally valid synthetic OD matrices via Dirichlet sampling, marginal perturbation, and IPF---scalable to $10^6$ samples. Part~II constructs turning-count observations through a route-choice-aware fractional operator $\Aturn$ ($\mathrm{rank}=174$, up from 122 under binary shortest-path assignment) on the Sioux Falls network, with optional Poisson noise and memory-mapped batch generation. Part~III specifies Ridge/Elastic Net baselines and MLP extensions for the underdetermined inverse problem, with regularization and forward-consistency losses motivated by null-space dimension $\mathcal{N}=402$. Future parts will add real-data validation, overdispersed noise, and equilibrium-based assignment models.

\begin{thebibliography}{9}
\bibitem{furness1962}
R.~Furness, \emph{Trip forecasting}, Traffic Engineering and Control, 1962.

\bibitem{deming1940}
W.~E.~Deming and F.~F.~Stephan,
\emph{On a least squares adjustment of a sampled frequency table when the expected marginal totals are known},
Annals of Mathematical Statistics, 1940.

\bibitem{ortuzar2011}
J.~de~D.~Ort\'{u}zar and L.~G.~Willumsen,
\emph{Modelling Transport}, 4th ed., Wiley, 2011.

\bibitem{cascetta1984}
E.~Cascetta,
\emph{Estimation of trip matrices from traffic counts and survey data},
Transportation Research Part B, 1984.
\end{thebibliography}

\end{document}
